% ECONOMICS_PROBLEM: {"id": "C78-352", "title": "Minimum reachable-state count for doubly prime stable-marriage instances", "jel_code": "C78", "source_id": "OP-0204"}
% !TeX program = xelatex
\documentclass[11pt,a4paper]{article}
\usepackage[margin=23mm]{geometry}
\usepackage{fontspec}
\setmainfont{Latin Modern Roman}
\usepackage{amsmath,amssymb,mathtools}
\usepackage{xcolor,enumitem,booktabs,longtable,array,xurl,hyperref}
\definecolor{muted}{HTML}{53636C}
\hypersetup{colorlinks=true,urlcolor=blue,linkcolor=blue}
\setlength{\parindent}{0pt}
\setlength{\parskip}{5pt}
\newcommand{\R}{\mathbb R}
\newcommand{\E}{\mathbb E}
\newcommand{\N}{\mathbb N}
\newcommand{\ind}{\mathbf 1}
\newcommand{\Var}{\operatorname{Var}}
\newcommand{\Cov}{\operatorname{Cov}}
\newcommand{\supp}{\operatorname{supp}}
\DeclareMathOperator*{\argmin}{arg\,min}
\DeclareMathOperator*{\argmax}{arg\,max}
\newcommand{\entrymeta}[3]{{\sffamily\small\color{muted}\textbf{\texttt{#1}}\quad|\quad #2\par #3\par}}
\newcommand{\Problemref}[1]{\texttt{#1}}
\newcommand{\JELref}[2]{\texttt{#2} (JEL \texttt{#1})}
\begin{document}
\section*{Minimum reachable-state count for doubly prime stable-marriage instances}
\textbf{Problem ID:} \texttt{C78-352}\quad \textbf{JEL:} \texttt{C78}\par
% BEGIN ECONOMICS STATEMENT
\label{problem:OP-0204}
% GAME_THEORY_PROBLEM: {"local_id": "GT-074", "original_number": 74, "kind": "P", "entry_kind": "statement_only", "published": false, "review_required": true, "category": "game_theory"}
\label{game:GT-074}
{\sffamily\small\color{muted}
\noindent\textbf{ID: GT-074.}\quad\textbf{Category:} Deferred acceptance combinatorics.
\quad\textbf{Kind: P.}\quad\textbf{Status:} Conjecture in the cited source;
larger reported minima are search bounds.
\par}
\subsection*{Definitions and mathematical statement}
An order-\(n\) instance \(I\) has disjoint proposer and receiver sets
\(P,R\), each of size \(n\), and strict complete rankings of the
opposite set. Write \(L_p(t)\) for proposer \(p\)'s \(t\)-th choice.
A deferred-acceptance state is \((h,q)\), where
\[
 h:R\to P\cup\{\bot\},\qquad q:P\to\{0,\ldots,n\},
\]
and the nonempty values of \(h\) are distinct. Initially \(h(r)=\bot\)
and \(q(p)=0\). A legal step chooses an unheld proposer \(p\)
with \(q(p)<n\), increments \(q(p)\), and proposes to \(L_p(q(p))\)
using the new value. That receiver keeps its preferred proposer among
its previous holder and \(p\), treating \(\bot\) as worse than every
proposer; the other proposer, if any, becomes unheld.
Let \(\operatorname{Reach}(I)\) contain every state obtainable by
finitely many legal steps, including the initial and terminal states.

An \emph{input-independent partition} consists of
\[
 P=P_1\sqcup\cdots\sqcup P_t,\quad
 R=R_1\sqcup\cdots\sqcup R_t,\quad |P_j|=|R_j|>0,
\]
such that each \(p\in P_j\) ranks all of \(R_j\) above every receiver
outside \(R_j\), and each \(r\in R_j\) ranks all of \(P_j\) above
every proposer outside \(P_j\).
The instance is \emph{input-prime} if no such partition has \(t\geq 2\).

For a reachable pointer vector \(q\), set
\[
 \begin{aligned}
 F(q)&=\{(p,t):p\in P,\ t\in\mathbb Z,\ 0\leq t<q(p)\},\\
 \mathcal F(I)&=\{F(q):(h,q)\in\operatorname{Reach}(I)\},\qquad
 E(I)=\bigcup_{X\in\mathcal F(I)}X.
 \end{aligned}
\]
The instance is \emph{behaviorally prime} if there is no partition
\(E(I)=E_1\sqcup E_2\) into two nonempty sets for which
\[
 \mathcal F(I)=
 \{X_1\cup X_2:
       X_j\in\{X\cap E_j:X\in\mathcal F(I)\},\ j=1,2\}.
\]
It is \emph{doubly prime} if it is both input-prime and behaviorally
prime. Let \(\mathcal D_n\) be the set of doubly prime order-\(n\)
instances, and set
\[
 m_n=\min_{I\in\mathcal D_n}|\operatorname{Reach}(I)|,
 \qquad \min\varnothing:=+\infty.
\]

\paragraph{Conjecture.}
For every integer \(n\geq 2\),
\[
                         m_n=3\cdot 2^{\,n-1}-1.
\]

\subsection*{Short English statement}
Is this formula the smallest possible number of execution states for
a doubly prime instance?

\subsection*{Variants and scope boundaries}
The count includes all serial schedules, not one run or only stable
matchings. Double primality is defined by the exact event-family product
condition. A computational graph test proposed in [S1] has a converse
that the paper reports as verified experimentally rather than proved;
it is not substituted for this definition.

\subsection*{Source relationship and status}
This is [S1, Section 7, problem A].
The proved lower bound is \(2^n+n-1\).
The values \(5,11\) are exhaustive minima at \(n=2,3\);
the values \(23,47,95\) at \(n=4,5,6\) are only upper bounds obtained
by constrained search. The catalogue overstated their exactness.

\subsection*{References}
\begingroup\footnotesize\raggedright
\noindent[S1] Yoshiteru Ishida,
\emph{Prime Factorisation of Stable-Matching Instances: Uniqueness,
Simultaneous Products, and an Exact Census}, 2026,
arXiv:2610.05617v1.
\url{https://arxiv.org/html/2610.05617v1}.
Locators: Section 2 (states and events);
Sections 3--4 (input and behavioral primality);
Proposition 21.4; Section 7, problem A.
\par\endgroup

\subsection*{Common conventions: Source mathematical conventions}

\textbf{Finite games.} For a finite set $A$, $\Delta(A)$ is its probability
simplex. Players choose independent mixed strategies in Nash equilibrium;
correlation is allowed only when explicitly part of the solution concept.
In a game with payoff functions $u_i$ and strategy profile $x$, an additive
$\varepsilon$ Nash equilibrium satisfies
\[
 u_i(x)\geq u_i(x_i',x_{-i})-\varepsilon
 \quad\text{for every player $i$ and every admissible deviation $x_i'$}.
\]
For numerical approximation, payoffs are normalized as locally specified.
An $\varepsilon$ payoff residual is distinct from distance at most
$\varepsilon$ to the set of exact equilibria. Point convergence, convergence
of empirical play, and convergence in distance to an equilibrium set are
also distinct.

\textbf{Computational representations.} Unless stated otherwise, a complexity
question takes rational data with binary encoding; polynomial time is measured
in total bit length. A fixed-accuracy polynomial algorithm is not necessarily
polynomial in $1/\varepsilon$, and neither is necessarily polynomial in
$\log(1/\varepsilon)$. Succinct games and value, demand, or sampling oracles
must declare their interfaces. Exact real solutions require an explicit
representation; an unspecified list of arbitrary real numbers is not a
finite computational output. A claimed algorithmic barrier is meaningful
only for its specified model and reductions.

\textbf{Dynamic games.} Histories, information, observation, and allowable
randomization are part of the model. A stationary strategy depends only on
the current state; a positional strategy in a deterministic graph game
chooses one action at each controlled vertex. Nash equilibrium at an initial
state and equilibrium after every history are different requirements.
An expectation of a pathwise $\liminf$ payoff, a $\liminf$ of expected averages,
a limiting expected average, and a uniform finite-horizon equilibrium are
different criteria. None is silently substituted for another.

\textbf{Allocation and mechanisms.} A complete allocation assigns every item
exactly once unless otherwise stated. Zero-valued goods, negative values,
capacity constraints, feasibility, lotteries, and copies of goods affect
fairness definitions and are specified locally. Dominant-strategy incentive
compatibility, Bayesian incentive compatibility, universal truthfulness,
and truthfulness in expectation impose different inequalities. Whenever
an objective ratio has zero denominator, its convention is stated or those
instances are excluded.

\textbf{Infinite-dimensional models.} Expectations, integrals, stochastic
equations, and optimization objectives require their local regularity,
integrability, filtrations, and admissible domains. A backward--forward
mean-field system is a boundary-value problem; it is not an initial-value
semiflow without an additional construction. Long-time, large-population,
and vanishing-noise limits require an order of limits and a convergence
topology. If a minimum need not be attained, the objective is its infimum
or supremum and the approximation requirement is stated separately.
% END ECONOMICS STATEMENT
\end{document}
